% !TEX root = ../main.tex
\chapter{Phenomenological Reference Atlas}
\label{app:paper_eqs_reference_atlas}

This appendix collects the archived paper equations figure material that was used as phenomenological reference during the transition to the QuTiP Lindblad and Redfield workflows. In the main text, these figures are no longer the primary results. They are retained here so that the reader can still compare the new QuTiP-based figures against the older phenomenological benchmark.

Some of the monomer archive figures are older thesis-local composites rather than freshly regenerated exports from one fully standardised run set. Where that is the case, they should be read as archived comparison material, not as the definitive baseline used in the main text.

\section{Monomer Archive}
\label{app:sec:paper_eqs_monomer}

\begin{figure}[htbp]
    \centering
    \safefallbackgraphic[0.92\linewidth]{figures/rslts_monomer_1d_hom_vs_inhom_fig2}{figures/rslts_monomer_1d_hom_vs_inhom_(fig2)}
    \caption{Archived monomer Fig.~2 comparison material retained from the phenomenological workflow. This figure is kept as a reference atlas entry rather than as part of the standardized Lindblad baseline in the main text.}
    \label{fig:app_paper_eqs_monomer_fig2}
\end{figure}

\begin{figure}[htbp]
    \centering
    \safefallbackgraphic[0.98\linewidth]{figures/rslts_monomer_2d_hom_fig3}{figures/rslts_monomer_2d_hom_(fig3)}
    \caption{Archived monomer Fig.~3 homogeneous 2D spectra from the phenomenological solver. This reference set is preserved so that the main-text Lindblad reproduction can always be compared back to the earlier archive.}
    \label{fig:app_paper_eqs_monomer_fig3}
\end{figure}

\begin{figure}[htbp]
    \centering
    \safefallbackgraphic[0.98\linewidth]{figures/rslts_monomer_2d_inhom_fig4}{figures/rslts_monomer_2d_inhom_(fig4)}
    \caption{Archived monomer Fig.~4 inhomogeneous 2D spectra from the phenomenological solver. As with~\cref{fig:app_paper_eqs_monomer_fig2}, this is retained as archived comparison material and not as the standardized main-text baseline.}
    \label{fig:app_paper_eqs_monomer_fig4}
\end{figure}

\section{Dimer Archive}
\label{app:sec:paper_eqs_dimer}

\begin{figure}[p]
    \centering
    \begin{minipage}[t]{0.49\linewidth}
        \centering
        \safeplaceholdergraphic[0.98\linewidth]{\dimerFigThreeUncoupledPaperPath}
    \end{minipage}\hfill
    \begin{minipage}[t]{0.49\linewidth}
        \centering
        \safeplaceholdergraphic[0.98\linewidth]{\dimerFigThreeCoupledPaperPath}
    \end{minipage}
    \caption{Archived dimer Fig.~3 reference set from the phenomenological solver: uncoupled control (left) and coupled dimer (right). This figure is retained as comparison material for the Lindblad calibration discussed in the main text.}
    \label{fig:app_paper_eqs_dimer_fig3}
\end{figure}

\begin{figure}[p]
    \centering
    \begin{minipage}[t]{0.49\linewidth}
        \centering
        \safeplaceholdergraphic[0.98\linewidth]{\dimerFigFourWaitFortySixPaperPath}
    \end{minipage}\hfill
    \begin{minipage}[t]{0.49\linewidth}
        \centering
        \safeplaceholdergraphic[0.98\linewidth]{\dimerFigFourWaitThreeTenPaperPath}
    \end{minipage}
    \caption{Archived dimer Fig.~4 waiting-time reference panels from the phenomenological solver at \(T=\SI{46}{\femto\second}\) and \(T=\SI{310}{\femto\second}\). They are kept here so that the Lindblad and Redfield dimer panels in the main text can always be compared back to the original phenomenological atlas.}
    \label{fig:app_paper_eqs_dimer_fig4}
\end{figure}

\begin{figure}[htbp]
    \centering
    \safeplaceholdergraphic[0.98\linewidth]{\dimerFigSixWaitSixtyTwoPaperPath}
    \caption{Archived dimer Fig.~6 homogeneous companion calculation from the phenomenological solver. This completes the archived dimer comparison atlas that accompanies the main-text Lindblad and Redfield discussion.}
    \label{fig:app_paper_eqs_dimer_fig6}
\end{figure}
